\documentclass[10pt]{article}
\usepackage[margin=20mm]{geometry}
\usepackage{amsmath,amssymb,mathtools,fontspec,unicode-math,xurl,hyperref}
\IfFontExistsTF{DejaVu Serif}{\setmainfont{DejaVu Serif}}{\setmainfont{Latin Modern Roman}}
\setmathfont{latinmodern-math.otf}
\setlength{\parindent}{0pt}\setlength{\parskip}{6pt}\sloppy
\title{Scope, spectral estimates and the good-cover argument for arithmetic locally symmetric spaces}
\author{AI-assisted technical review}
\date{}
\begin{document}
\maketitle
Abstract: The acquired revised manuscript of Fraczyk, Hurtado and Raimbault is examined from its arithmetic collar lemma through orbital growth, thin-volume estimates and the bounded-degree nerve construction. The introductory collar theorem lacks an essential irreducibility restriction. The rank-one spectral-gap calculation and the final nerve-degree argument contain explicit errors in core derivations. Several narrower algebraic and tube-counting statements also require correction. Counterexamples, exact checks and local replacement arguments are supplied, including a factorwise supplement for the main theorem. The overall assessment is Major for the specified claims and derivations. These findings do not establish that the corrected topological-complexity or torsion-homology conclusions are false.\par
\section*{Identity, exact version and mathematical scope}
The reviewed paper is Mikołaj Frączyk, Sebastian Hurtado and Jean Raimbault, Topological complexity of arithmetic locally symmetric spaces, associated with Inventiones mathematicae (2025), DOI 10.1007/s00222-025-01394-1 (https://doi.org/10.1007/s00222-025-01394-1). The actual source is arXiv:2202.13619v2, dated 2 February 2026, 23 pages [1]. No equivalence with the publisher’s version is asserted. References to the source below use physical PDF pages, equal to its printed page numbers. The accompanying `mainrev.tex`, `macros.tex` and bibliography were read as source text; they were not executed.\par
For a fixed symmetric space \(X\), Theorem A asserts that every closed arithmetic locally \(X\)-manifold has a homotopy-equivalent simplicial complex with both its total number of simplices and its vertex incidences controlled: the former is linear in volume, the latter uniform in the manifold. Theorem B bounds the logarithm of integral homology torsion for finite-volume arithmetic manifolds. The noncompact topological case is an earlier Gelander input, not newly proved in this manuscript. The new compact argument passes through the degree-dependent arithmetic Margulis lemma, orbital growth, a polynomial overlap estimate and an adaptive good cover.\par
Write \(d=[k:\mathbb Q]\), \(n=\dim X\), and use the ordinary geometric convention \[\operatorname{inj}_M(x)=\frac12\inf_{\gamma\ne1}d_X(\widetilde x,\gamma\widetilde x)\] for a torsion-free effective lattice quotient. For orbifolds this expression measures proximity to the singular locus. The source’s first injectivity-radius display on p. 8 has a factor inconsistent with its following \(B(2R)\) display; its intended standard convention is used here. Accordingly tubes defined by displacement \(<R\) correspond to injectivity radius \(<R/2\). Constant rescaling reconciles the later volume estimates.\par
All six principal sections and the exact author TeX were read in this initial review; the primary session actually inspected source images pp. 1–21. Selected passages of the primary dependencies [2]–[6] were also read and viewed. Their complete proofs, the full arithmetic-lattice classification, the flat torus theorem, Jordan’s theorem and the quoted uniform triangularization results remain named external inputs.\par
\section*{Problem 1: the introductory collar theorem needs irreducibility}
\subsection*{The quantifier failure}
Theorem C [1, p. 3] says that for any arithmetic lattice in the fixed group, elements moving a point by at most \(\varepsilon_X d\) generate a virtually nilpotent subgroup. It does not impose irreducibility. The actual Theorem 3.1, p. 9, does impose it.\par
Take \(G=\operatorname{PSL}_2(\mathbb R)\times\operatorname{PSL}_2(\mathbb R)\). In its first factor use the arithmetic lattice \(\Gamma_1=\Gamma(2)/\{\pm I\}\), containing \[a=\begin{pmatrix}1&2\\0&1\end{pmatrix},
 \qquad b=\begin{pmatrix}1&0\\2&1\end{pmatrix}.\] Their nonzero powers act by ping-pong on the projective-line regions \(\{|x|<1\}\) and \(\{|x|>1\}\), respectively; hence \(\langle a,b\rangle\) is a nonabelian free group. In the curvature \(-1\) normalization, both move \(i\in\mathbb H^2\) by \(\operatorname{arcosh}(3)\), since \[\cosh d(i,gi)=\frac{a_g^2+b_g^2+c_g^2+d_g^2}{2}.\] A different fixed normalization of \(X\) merely changes this fixed displacement.\par
Choose totally real fields \(F_j\) of unbounded degree and quaternion algebras over \(F_j\) split at exactly one real place. The standard norm-one order construction gives arithmetic Fuchsian lattices \(\Gamma_j'\) with invariant, equivalently adjoint, trace field \(F_j\); the construction and field identification are explicitly recorded in [2, p. 534]. For degree greater than one these may be chosen cocompact. Set \(\Gamma_j=\Gamma_1\times\Gamma_j'\). This is an arithmetic lattice in the same fixed \(G\). Its adjoint trace field is \(F_j\): traces on the direct-sum Lie algebra add, the first factor has rational adjoint traces, and \(\operatorname{tr}\operatorname{Ad}(1,\gamma)=3+\operatorname{tr}\operatorname{Ad}(\gamma)\).\par
For every proposed \(\varepsilon_X>0\), sufficiently large \(j\) makes \(\varepsilon_X[F_j:\mathbb Q]\) exceed the fixed displacement of \(a\) and \(b\). The short-element subgroup at \((i,x_j)\) then contains \(\langle(a,1),(b,1)\rangle\cong F_2\). A virtually nilpotent group cannot contain this subgroup. Every condition in Theorem C as printed is met; its conclusion fails.\par
\subsection*{Repair and impact}
Add irreducible to Theorem C, as already required by Theorem 3.1, or formulate a factorwise radius using the individual irreducible degrees. A single full product trace degree cannot replace these separate parameters. This is Major for a prominently stated technical guarantee central to the paper’s method. It does not refute Theorem 3.1. The proof of Theorem A on p. 18 also applies irreducible degree estimates to an unrestricted lattice without explaining that reduction; a valid factorwise completion is given in Section 8 (\#sec:completion) below.\par
For the irreducible compact case, the height-gap mechanism itself checks out after the local height correction in Problem 3. The small-displacement set is finite, contains the identity, and its word displacements are bounded by word length times the original bound. Breuillard’s actual Corollary 1.7 [3, p. 1061] supplies the stated uniformly short word with an eigenvalue of positive height. The anisotropic solvable-subgroup argument then gives virtual abelianness. The nonuniform rational-structure sentence should be read with its cited Gelander Lemma 5.2 [4, physical pp. 20–21]; a naive objection based only on a presentation over a larger number field was rejected after that passage was checked.\par
\section*{Problem 2: the orbital-growth proof uses an incorrect spectral calculation}
\subsection*{A decisive rank-one test}
Lemma 4.2 [1, p. 11] assigns critical exponents \(m,m,2m\) to its three listed real and complex hyperbolic subgroup cases. Its spectral formula uses \(\dim\partial X\) as the entropy parameter. Already the first case is inconsistent with the desired conclusion. Take \(G=\operatorname{SO}^0(2,1)\), \(H=\operatorname{SO}^0(1,1)\), the centralizer of a hyperbolic element, and a cyclic uniform lattice \(\Lambda\) in \(H\). The orbit of \(\Lambda\) grows linearly along a geodesic, so its critical exponent is zero. The printed assignment has \(m=1\), hence critical exponent one. Together with \(\dim\partial\mathbb H^2=1\), the displayed second branch gives \(\lambda_0=1(1-1)=0\), not a positive lower bound. This is an admissible centralizer, not an excluded endpoint.\par
The correct metric-invariant parameter is the volume entropy \(h_X=2\|\rho\|\). The rank-one formula, explicitly stated in [5, p. 920], is \[\lambda_0(\Lambda\backslash X)=
 \begin{cases}
 h_X^2/4,&0\le\delta_\Lambda\le h_X/2,\\
 \delta_\Lambda(h_X-\delta_\Lambda),&h_X/2\le\delta_\Lambda\le h_X .
 \end{cases}
 \tag{3.1}\] For real hyperbolic \(m\)-space of curvature \(-1\), the lattice exponent is \(m-1\), with the cyclic \(m=1\) case giving zero. For complex hyperbolic \(n\)-space in the usual holomorphic-curvature \(-4\) normalization, \(h_X=2n\), whereas its topological boundary has dimension \(2n-1\). The real and complex subgroup exponents in that normalization are \(m-1\) and \(2m\). Metric rescaling must be applied consistently to all of these numbers; it cannot turn linear orbit growth into positive exponential growth.\par
These corrected values are strictly below \(h_X\) for every listed proper noncompact subgroup, with only finitely many relevant dimensions for fixed \(G\). Compact centralizers are handled by the regular-representation spectral bound. Thus the corrected formula supplies the intended positive uniform gap. The property-\((T)\) branch is separate; one can choose the averaging kernel as a convolution square to ensure positivity of the self-adjoint averaging operator.\par
\subsection*{The compact-factor restriction}
Theorems E and 4.1 [1, pp. 4, 10] also state their exponential orbital-growth assertion for arbitrary semisimple \(G\) with finite center. Lemma 4.2 excludes compact factors. This discrepancy has a literal counterexample. Let \[G=G_0\times\operatorname{SO}(3),\qquad
 \Gamma=\Gamma_0\times F,
 \qquad \gamma=(1,u),\] where \(\Gamma_0\) is a uniform lattice in the noncompact \(G_0\), \(F\) is the finite tetrahedral rotation group, and \(u\) is an order-three rotation noncentral in both \(F\) and \(\operatorname{SO}(3)\). Then \(\Gamma\) is uniform and \(\gamma\) is noncentral. The compact factor acts trivially on \(X=G/K\), so every conjugate of \(\gamma\) has displacement zero. Its orbital integral over \(G_\gamma\backslash G\) is the same finite positive compact-orbit volume for every positive radius. The asserted inequality would force \(1<C e^{-\delta R_2}\) for arbitrarily large \(R_2\), which is impossible.\par
Restrict the theorem to groups without compact factors, or require a noncentral projection to the effective noncompact quotient. The earlier no-compact-factor phrase on p. 5 is explicitly local to its paragraph; it does not repair the independently restated broader theorems. Replace the strict inequality by a non-strict one at zero orbital mass.\par
\subsection*{Core dependency and counterevidence}
The spectrum calculation is Major as an invalid derivation of the uniform gap used in \[\mathrm{Lemma}\ 4.2\ \Longrightarrow\ \mathrm{Theorem}\ 4.1\
 \Longrightarrow\ \mathrm{Theorem}\ 5.7\ \Longrightarrow\ \mathrm{Theorem}\ A.\] The compact-factor counterexample also refutes the literal orbital-growth guarantee. Neither finding disproves its corrected effective version. Once a genuine gap is supplied, the p. 12 convolution argument is sound: convolution expands displacement support from \(R\) to \(R+2\), preserves the integral, and Cauchy–Schwarz gives \[\mathcal O(\gamma,1_{B(R)})\le e^{-2\delta}\mathcal O(\gamma,1_{B(R+2)}).\] Iteration and monotonicity yield the required exponential growth, with a uniform prefactor.\par
\section*{Problem 3: the auxiliary splitting-field bound is false}
In the proof of Proposition 2.2 [1, p. 8], the manuscript asserts \[[k(\Phi(\gamma)):k]\le n,\qquad n=\dim\mathbf G,\] where \(k(\Phi(\gamma))\) is generated by all adjoint root values. Individual eigenvalue degrees have a dimension bound, but their common splitting field need not.\par
Take \(G=\operatorname{PGL}_5(\mathbb R)\), the arithmetic lattice given by the image of \(\operatorname{SL}_5(\mathbb Z)\), \(k=\mathbb Q\), and the companion matrix of \(f(t)=t^5-t-1\). It has determinant one and distinct roots, so its image is an allowed semisimple element. Exact polynomial division gives irreducibility modulo three and the square-free factorization \[f(t)\equiv(t^2+t+1)(t^3+t^2+1)\pmod2.\] Thus the Galois group is transitive of prime degree five and contains a permutation of cycle type \((2)(3)\), whose cube is a transposition. A primitive group containing a transposition is the full symmetric group: the graph of its conjugate transpositions has components forming a block system. Consequently the splitting field has degree \(120\).\par
Its root ratios generate that whole field. An automorphism fixing every ratio scales every root by a common \(\zeta\); the nonzero fourth and fifth elementary symmetric functions force \(\zeta^4=\zeta^5=1\), hence \(\zeta=1\). Therefore \[[\mathbb Q(\Phi(\gamma)):\mathbb Q]=120>24=\dim\operatorname{PGL}_5.\] The exact modular tests are preserved in `aux\_height/exact\_checks.json`; the Galois and ratio-field arguments are the proof, not a numerical extrapolation.\par
There is a direct repair requiring no splitting-field bound. Let \(L=k(\Phi(\gamma))\). For an archimedean embedding \(\sigma\) of \(k\), set \[S_\sigma=\sum_{\alpha\in\Phi}\log^+|\tau(\alpha(\gamma))|,\] where \(\tau\) is any extension to \(L\). Galois invariance of the root-value multiset makes \(S_\sigma\) independent of that choice. Each \(\sigma\) has exactly \([L:k]\) extensions; this multiplicity cancels the denominator in Weil height. If one representative is used per complex conjugate pair, the exact relation is \[[k:\mathbb Q]h(\gamma)=
 \sum_{\sigma\ \mathrm{real}}S_\sigma+
 2\sum_{\sigma\ \mathrm{complex\ pairs}}S_\sigma.
 \tag{4.1}\] Thus, writing \(S\) for the corresponding unweighted representative sum, \(S\le[k:\mathbb Q]h(\gamma)\le2S\); the p. 8 displayed inequality has the wrong direction when complex-place contributions occur. Together with the preceding displacement comparison, (4.1) restores Proposition 2.2 with constants depending only on \(X\). For zero height, state it multiplicatively instead of using an undefined \(0/0\) quotient. The false extra field-degree bound is Minor: it is unnecessary, and the already available multiset invariance supplies the full intended comparison.\par
\section*{Problem 4: the overlap bound loses its necessary uniform prefactor}
Proposition 5.1 [1, p. 13] bounds the number of small-displacement tubes through a point by \(d^m\), with no multiplicative constant. At \(d=1\) this is exactly one, independent of \(m\).\par
Consider \[q=x_1^2+x_2^2+x_3^2-7x_4^2,\qquad
 \Gamma=\operatorname{SO}^+(q,\mathbb Z).\] The form is anisotropic over \(\mathbb Q\). In a primitive integral solution of \(q=0\), an odd \(x_4\) would express seven modulo eight as a sum of three squares; an even \(x_4\) forces all three other coordinates even, contradicting primitivity by descent. The standard anisotropic arithmetic-lattice criterion therefore gives a uniform lattice in the real simple group of signature \((3,1)\). Its adjoint trace field is \(\mathbb Q\).\par
The integral rotation \(r\) through \(90^\circ\) in the first two coordinates preserves \(q\), has order four, and fixes a geodesic in hyperbolic three-space. Both \(r\) and \(r^2\) have zero displacement on that axis. Their positive-radius tubes are genuinely distinct: in curvature \(-1\), \[\sinh\!\left(\frac{d(x,r_\theta x)}2\right)
 =|\sin(\theta/2)|\sinh d(x,\mathrm{axis}),\] so \(S_{r^2}^{R}\subsetneq S_r^{R}\). Thus every positive radius gives at least two distinct tubes through an axis point, contradicting \(d^m=1\). The exact residue and matrix checks are saved in `aux\_thin/exact\_checks.json`. The criterion establishing arithmetic cocompactness is a named standard input, not inferred from finite enumeration.\par
Replace the bound by \(C(d+1)^m\), or \(Cd^m\) with the small-degree cases treated explicitly. This is a Minor boundary defect in the stated overlap bound; its proof itself already produces multiplicative constants, and the required large-degree polynomial scale is unchanged.\par
Several supporting proof steps on pp. 13–16 also need precise completion. Lemma 5.2’s hypothesis implies that every element is semisimple: a positive power belongs to its semisimple abelian finite-index subgroup, and in characteristic zero this forces the original unipotent Jordan part to vanish. Hence the alternative uniform triangularization argument already given in Remark 5.3 applies, conditional on its stated external inputs. It avoids the unjustified identification of the chosen torsion-free subgroup with the entire diagonal intersection.\par
For the flat action, normalizing a translation subgroup does not preserve an arbitrary flat selected from its minimum set. Choose a normal abelian subgroup of bounded index, then a characteristic torsion-free finite-index subgroup \(L_0\). The transverse factor in \(\operatorname{Min}(L_0)=Y\times E\) has a fixed point under the finite quotient. This selects a \(\Delta\)-invariant flat. Its action kernel is finite, and Jordan’s theorem together with bounds on algebraic roots of unity bounds that kernel polynomially in \(d\). The useful elementary bound is \(\varphi(N)\ge\sqrt{N/2}\), not the printed \(\sqrt N\), which fails for \(N=2,6\).\par
The translation lattice of the projected group need not be exactly \(\pi(L_0)\). For any projected nonzero translation, equivariant nearest-point projection to the chosen flat shows that its translation length equals that of its lift on \(X\). The nonelliptic systole bound therefore applies directly. At bounded degree use a positive constant, or write the coarse bound with \(\log(d+2)\). The rank-zero case is just a finite-group bound.\par
Finally, the successive-minima vectors in Lemma 5.5 need not form a lattice basis. For \[L=\mathbb Z^4+\mathbb Z(1/2,1/2,1/2,1/2)\] the four standard unit vectors attain all successive minima but generate an index-two sublattice. The printed fundamental-domain assertion therefore fails. A direct replacement is simpler: in each coset of the translation lattice, orbit points are separated by its minimum length \(\ell\). Disjoint balls of radius \(\ell/2\) give \[\#\{g:d(gx,x)\le R\}\le A(r)(1+2R/\ell)^r,
 \qquad r=\dim E,
 \tag{5.1}\] where the finite point-group bound \(A(r)\) depends only on \(r\). This covers \(R<\ell\) as well. Together these completions recover the intended polynomial overlap estimate, with explicit external Jordan, flat-torus and triangularization dependencies.\par
\section*{Problem 5: the tube-volume comparison is false in the nonuniform range}
Corollary 5.6 [1, p. 16] omits the uniformity assumption of Proposition 5.1. It claims that a finite polynomial in \(d\) controls the sum of volumes of all conjugacy-class tubes by the thin-part volume. No finite constant can make that assertion true for every nonuniform arithmetic lattice.\par
Take \(\Gamma=\operatorname{PSL}_2(\mathbb Z)\), \(d=1\), and \(\tau^j(z)=z+j\) for positive integers \(j\). The centralizer in \(\Gamma\) is generated by \(\tau\). The displacement formula is \[d(iy,iy+j)=2\operatorname{arsinh}\frac{j}{2y}.\] Thus the displacement-\(<R\) tube has, in the centralizer quotient, area \[\int_0^1\int_{j/(2\sinh(R/2))}^{\infty}\frac{dy\,dx}{y^2}
 =\frac{2\sinh(R/2)}{j}.
 \tag{6.1}\] With \(\operatorname{vol}K=1\), this is the corresponding group-quotient tube volume. The positive powers represent distinct conjugacy classes. Indeed, the gcd of the entries of \(\tau^j-I\) is \(j\) and is invariant under conjugation by integral unimodular matrices; the sign ambiguity in \(\operatorname{PSL}_2\) is removed by taking the trace-two representative.\par
For every \(R>0\), including every sufficiently small allowed radius, (6.1) yields a divergent harmonic class sum. The modular surface has finite area, so its thin part has finite area. The asserted lower comparison therefore fails.\par
Restrict Corollary 5.6 to uniform lattices and retain the prefactor from Problem 4. For a rigorous unfolding, count fibers over \(\Gamma\backslash G\), normalize Haar measure on \(K\), and then descend to \(\Gamma\backslash X\). The literal map displayed to \(X\) has continuous \(K\)-fibers and is not finite-to-one. Also use displacement \(2R\) when the thin radius is defined by the usual injectivity radius.\par
This is Minor in its established impact: the false nonuniform extension is auxiliary to the new compact-manifold argument, which uses uniform lattices. The bounded-degree nonuniform range of Theorem 5.7 can be obtained separately from monotonicity by increasing its constant. That observation does not prove the strictly smaller-than-one fraction asserted in Theorem D; this distinction is retained below.\par
\section*{Problem 6: the good-cover proof does not count the adjacent vertices}
\subsection*{The essential wrong set}
In the final paragraph proving Proposition 6.1 [1, p. 20], the degree bound counts centers \(q\in B(p,\varepsilon_p/6)\). Adjacency in the nerve instead means \[B(p,\varepsilon_p/6)\cap B(q,\varepsilon_q/6)\ne\varnothing.\] The latter does not require that \(q\) lie in the first ball. Equal radii \(1/6\) and center distance \(1/4\) give a direct local example. The two centers also satisfy the required net separation. Thus the displayed count omits allowed neighbors and does not establish the claimed vertex degree. This is the final essential step supplying the bounded-incidence conclusion of Proposition 6.1 and Theorem A.\par
There is a complete replacement. The radius function \(\varepsilon_x=\min(1,\operatorname{inj}_M(x))\) is 1-Lipschitz. Writing \(s=d(p,q)\), intersecting balls imply \[6s<\varepsilon_p+\varepsilon_q
 \le2\varepsilon_p+s.\] Hence \[s<2\varepsilon_p/5,\qquad
 \varepsilon_q>\varepsilon_p-s>3\varepsilon_p/5.
 \tag{7.1}\] All neighbor centers therefore lie in \(B(p,2\varepsilon_p/5)\) and, by the net condition, are \(\varepsilon_p/60\)-separated. The corrected packing bound below gives \[\deg(p)\le C_X(49).\] Every simplex incident to \(p\) is determined by a subset of its neighbors, so the number of incident simplices is at most \(2^{\lceil C_X(49)\rceil}\). This restores precisely the missing uniform bound.\par
\subsection*{Two quantitative corrections needed in the same construction}
Lemma 6.2, p. 18, claims \(|P|<C_X(\delta/\eta)\) for an \(\eta\)-separated set in a \(\delta\)-ball. The correct packing argument lifts the centers to \(X\); their disjoint \(\eta/2\)-balls fit in the ball of radius \(\delta+\eta/2\). It gives \[|P|\le C_X(2\delta/\eta+1).
 \tag{7.2}\] The stated smaller constant is actually false. In curvature-\(-1\) \(\mathbb H^2\), choose \(\eta=1/10\), \(\delta=11/100\) and points at distances \(3/40\) on opposite sides of a geodesic midpoint. There are two separated points in the required ball, but \[C_{\mathbb H^2}(11/10)\le(3520/3079)^2<2.\] This rational estimate follows from \(\sinh(ru)/\sinh u\le r\cosh(ru)\) for \(0<u<1/4\), together with \(\cosh(11/40)\le(1-121/3200)^{-1}\). The exact arithmetic and the repaired constants are recorded in `aux\_nerve/exact\_constant\_checks.json`. For any fixed rescaling of the metric, continuity as \(\delta/\eta\) tends to one gives the same failure.\par
The function \(F_p\) in Lemma 6.3, pp. 18–19, has support radius \(\varepsilon_p/4\) but normalization denominator \(\operatorname{vol}B(p,\varepsilon_p/2)\). Consequently \[\int_M F_p=
 \frac{\operatorname{vol}B(p,\varepsilon_p/4)}{\operatorname{vol}B(p,\varepsilon_p/2)}
 <1,\] contrary to the printed equality. Use instead \(\int F_p\ge C_X(2)^{-1}\). Formula (7.2) changes the multiplicity bound for \(H_p(x)\) to \(C_X(73)\). The argument then gives \[|S|\le C_X(2)^2C_X(73)
 \int_M\operatorname{vol}B(x,\varepsilon_x/2)^{-1}\,dx
 \ll_X\int_M\varepsilon_x^{-n}\,dx.
 \tag{7.3}\] The embedded-ball volume is comparable to \(\varepsilon_x^n\), with its inverse comparable to \(\varepsilon_x^{-n}\).\par
\subsection*{What the construction does establish}
Maximality of the separated net does give coverage: using the Lipschitz property in the failed separation inequality yields \(d(x,p)\le\varepsilon_p/35<\varepsilon_p/6\). Compactness ensures a finite net. The source’s lift argument rules out distinct lifts intersecting the same small ball, because their center separation would be below the injectivity bound. All finite intersections then lift to convex ball intersections in \(X\), so the cover is good. The nerve theorem supplies the homotopy equivalence.\par
The assessment is Major for the wrong-set degree argument in a core construction, with the normalization and packing corrections as additional local defects. The replacement (7.1)–(7.3) verifies the intended bounded-degree nerve conclusion. A short and successful repair is not evidence that the original derivation was valid.\par
\section*{Completion of the main reduction and remaining dependencies}
\subsection*{Irreducible large and small degrees}
After the corrected uniform overlap and orbital-growth arguments, the thin-volume comparison gives \[\frac{\operatorname{vol}(M_{\le\eta d})}{\operatorname{vol}M}\le C d^m e^{-ad}.\] For sufficiently large \(d\), this is at most \(e^{-a'd}\) for a fixed \(0<a'<a\). Ensure also \(1\le\eta d\) before using it to estimate the 1-thin part. Combining it with the nonelliptic length bound gives \[e^{-a'd}\log(d+2)^{3n}=O_{X,n}(1).\] For bounded degree, the eigenvalue degrees are uniformly bounded and non-roots of unity have a positive uniform height lower bound; equivalently use the bounded-degree part of the systole argument. Thus the compact torsion-free case has a uniform positive injectivity lower bound in that range. No statement that a positive-radius thin estimate is vacuous is needed.\par
\subsection*{The reducible case, including projected torsion}
Use the standard decomposition of a lattice into discrete irreducible projection blocks. For the connected effective group, write \(\Gamma\subset\prod_i\Gamma_i\) with finite index, allowing torsion in the cocompact projected lattices \(\Gamma_i\). No uniform bound on that index is assumed.\par
For \(M_i=\Gamma_i\backslash X_i\) define the nonelliptic radius \[\rho_i(x_i)=
 \min\!\left(1,\frac12
 \inf_{\gamma_i\ \mathrm{of\ infinite\ order}}
 d_{X_i}(\widetilde x_i,\gamma_i\widetilde x_i)\right).
 \tag{8.1}\] Take the empty infimum as \(+\infty\). Every nonidentity element of the torsion-free \(\Gamma\) has an infinite-order projection: a finite tuple of finite-order elements has finite order. For the product metric, its displacement is at least that in this component. It follows that \[\varepsilon_M(x)\ge\min_i\rho_i(x_i),
 \qquad
 \varepsilon_M(x)^{-n}\le\sum_i\rho_i(x_i)^{-n}.
 \tag{8.2}\] Because each factor function descends to its orbifold quotient, finite-cover and product integration give exactly \[\frac1{\operatorname{vol}M}\int_M\rho_i(x_i)^{-n}\,dx
 =
 \frac1{\operatorname{vol}M_i}\int_{M_i}\rho_i(x_i)^{-n}\,dx_i.
 \tag{8.3}\] The covering index cancels.\par
At bounded factor degree, the positive nonelliptic systole bounds (8.1) away from zero, including near orbifold singularities. At large degree, \(\{\rho_i<1\}\) is contained in the usual 1-thin part, so the large-degree exponential bound and the logarithmic systole bound give a uniform upper bound for the right side of (8.3). Equations (8.2)–(8.3) and the corrected Proposition 6.1 therefore supply the general compact conclusion, conditional on the named standard lattice decomposition. This handles projected torsion; blindly using the ordinary orbifold injectivity radius would instead introduce zeros at singular strata.\par
If disconnected effective isometries are allowed, put \(\Gamma_0=\Gamma\cap\operatorname{Isom}(X)^0\) and \(f=[\Gamma:\Gamma_0]\). The component count bounds \(f\) solely in terms of \(X\). For every \(\gamma\ne1\), torsion-freeness gives \(\gamma^f\ne1\) in \(\Gamma_0\), and \[d(x,\gamma^f x)\le f\,d(x,\gamma x).\] Hence \(\varepsilon_\Gamma\ge\varepsilon_{\Gamma_0}/f\). Normalized integration over the \(f\)-sheeted cover increases the bound by at most \(f^n\), still uniform in \(M\). For connected effective groups the discrete irreducible projection decomposition itself needs no preliminary cover. These reductions replace the unjustified use of one product trace degree; they do not assume a uniform index in the product of the projected lattices.\par
\subsection*{Torsion and other claims}
For the compact case of Theorem B, the repaired nerve has \(O_X(\operatorname{vol}M)\) simplices. Each simplicial boundary column in degree \(i+1\) has Euclidean norm \(\sqrt{i+2}\). Hadamard’s determinant bound controls the torsion of its cokernel exponentially in the number of columns. The kernel of the next boundary is primitive, so this is also the relevant homology torsion bound. The noncompact case retains the previous Gelander topological result as a dependency.\par
Two ranges remain explicitly unresolved in this bounded review. First, §5.1, p. 17, does not establish the full bounded-degree nonuniform version of Theorem D: absorbing \(Cd^m\) works for large degree, and bounded degree alone does not make its strictly subunit volume fraction automatic. Monotonicity repairs Theorem 5.7 with a prefactor but does not supply this stronger assertion. This is recorded as Needs Evidence, without claiming a counterexample to Theorem D.\par
Second, the intended lower bound in Proposition 6.4 concerns \(C\simeq M\), as its proof states; the printed \(C\simeq X\) is an isolated target-symbol error. The consulted Lafont–Schmidt straightening [6, pp. 133, 135–137] has the required face compatibility and compatible equivariant homotopies. These justify gluing and preservation of the fundamental class within its scope. Its bounded-Jacobian theorem, however, excludes every local \(\mathbb H^2\) or \(\operatorname{SL}_3(\mathbb R)/\operatorname{SO}(3)\) factor. The present manuscript explicitly discusses the exceptional whole spaces but does not supply the mixed-factor bounded-image argument. A complete exceptional/product extension was not verified here; the bounded attempts to acquire the cited Bucher text failed and are recorded. This dependency remains Needs Evidence, not an asserted false lower-bound theorem. The remarks about adaptive Delaunay triangulation also retain their external construction as a dependency; no triangulation algorithm was executed.\par
\section*{Impact and review state}
The unsupported generality of the introductory collar theorem and orbital-growth theorem is separated from their corrected effective, irreducible uses. The spectral-gap and nerve-degree calculations are demonstrated invalid core derivations. The exact counterexamples to the height estimate, unit-prefactor overlap bound and nonuniform tube comparison have narrower impact. The replacement arguments above support the main compact-complexity chain without pretending that every cited external proof, every version of Theorem D, or the mixed-factor lower-bound extension has been independently reconstructed.\par
Concern | Grade | Repair verification | Adoption\par
Unrestricted collar radius | Major | Product counterexample; irreducible restriction and factorwise supplement | Unknown\par
Orbital growth and spectrum | Major | Effective-group scope and entropy formula checked | Unknown\par
Splitting-field estimate | Minor | Exact degree witness; height identity restored | Unknown\par
Overlap prefactor | Minor | Orbifold witness; polynomial bound completed | Unknown\par
Nonuniform tube sum | Minor | Divergent class sum; uniform restriction verified | Unknown\par
Good-cover degree | Major | Correct neighbor set, packing and mass bounds derived | Unknown\par
Two remaining extensions | Needs Evidence | Full nonuniform thin fraction and exceptional mixed-factor straightening unverified | Unknown\par
No author adoption or correction was independently established. Scientific adjudication remains pending, with `human\_reviews=[]`. This is an initial recovery review, not independent acceptance.\par
9 M. Frączyk, S. Hurtado and J. Raimbault, Topological complexity of arithmetic locally symmetric spaces, arXiv:2202.13619v2, 2 February 2026, 23 pp., https://arxiv.org/pdf/2202.13619v2. D. D. Long and A. W. Reid, On Fuchsian groups with the same set of axes, Bull. London Math. Soc. 30 (1998), 533–538. Consulted author-hosted PDF, physical p. 2/printed p. 534, https://math.rice.edu/\textasciitilde{}ar99/axes.pdf. E. Breuillard, A height gap theorem for finite subsets of GL(d, Qbar) and nonamenable subgroups, Ann. of Math. 174 (2011), 1057–1110. Consulted publisher PDF, Corollary 1.7, p. 1061. T. Gelander, Homotopy type and volume of locally symmetric manifolds, acquired arXiv:math/0111165v5; associated with Duke Math. J. 124 (2004), 459–515. Consulted physical pp. 20–21, especially Lemma 5.2. E. Leuzinger, Critical exponents of discrete groups and L2-spectrum, Proc. Amer. Math. Soc. 132 (2004), 919–927. Consulted AMS PDF, pp. 919–921, especially the rank-one formula on p. 920. J.-F. Lafont and B. Schmidt, Simplicial volume of closed locally symmetric spaces of non-compact type, Acta Math. 197 (2006), 129–143. Consulted author-hosted PDF, pp. 130, 132–133, 135–137, https://people.math.osu.edu/lafont.1/LafSch.pdf.\par
\section*{Supporting checks}
The exact canonical PDF has SHA256\par
a5961f671e007bc229f0b822896b57d4
d6a6bd4876569541001c27522460896e\par
with the two lines concatenated.\par
and the recovery token is\par
No substantive prior report existed in the assigned preserved snapshot. The existing acquired PDF, source manifest and author source files were reused after actual hash verification; inventory parsing was not counted as scientific inspection.\par
The primary session read all 673 lines of the exact main TeX, all 146 macro lines and all 197 bibliography lines, and actually viewed canonical pp. 1–21. Native bounded auxiliaries `/root/p0128\_height`, `/root/p0128\_thin` and `/root/p0128\_nerve` supplied separately saved arguments, exact computations and image observations. The primary integrated their evidence and checked the decisive counterexamples and replacements. The external controller’s Personal/Chat/GPT-6/Pro observation is retained as such; backend model identity was not independently verified.\par
Reviewer-written exact checks include modular polynomial division, quaternion-free quadratic-form and rotation witnesses, lattice-index and root-of-unity bounds, and rational good-cover constants. These computations support the explicitly stated proofs and do not substitute for the external structural theorems. No author code was executed.\par
The evidence records bind exact source and auxiliary PDFs, scripts, runtimes, renderers, logs and actually inspected images. The current self-contained report is compiled with the existing approved Tectonic executable using `--untrusted` and `--only-cached`. Compiler and complete isolated support-cache inventories are hashed before and after execution. Every report page is actually viewed before delivery completion; rendering and hashing alone are not visual inspection.\par
\end{document}
